\documentclass[11pt,a4paper]{article}
\usepackage[utf8]{inputenc}
\usepackage[T1]{fontenc}
\usepackage[ngerman]{babel}
\usepackage[margin=2cm]{geometry}
\usepackage{enumitem}
\usepackage{xcolor}
\usepackage{titlesec}
\usepackage{fancyhdr}
\usepackage{siunitx}

% Page style
\pagestyle{fancy}
\fancyhf{}
\fancyhead[L]{\small PhD-Verteidigung - Stichpunkte}
\fancyhead[R]{\small Seite \thepage}
\renewcommand{\headrulewidth}{0.4pt}

% Section formatting
\titleformat{\section}{\Large\bfseries\color{blue!70!black}}{\thesection}{1em}{}
\titleformat{\subsection}{\large\bfseries\color{blue!50!black}}{\thesubsection}{1em}{}
\titleformat{\subsubsection}{\normalsize\bfseries}{\thesubsubsection}{1em}{}

\definecolor{ipp}{rgb}{0,0.40,0.75}
\definecolor{ippalt}{rgb}{0.09,0.56,1.0}
\definecolor{ipplight}{rgb}{0.78,0.85,0.95}
\definecolor{ippdark}{rgb}{0,0.22,0.46}

% Compact lists
\setlist[itemize]{leftmargin=*,topsep=2pt,itemsep=1pt,parsep=0pt}

% Title
\title{\textbf{PhD-Verteidigung Präsentation}\\[0.5em]
\Large Bolometrie-Diagnostik und Echtzeit-Strahlungsrückkopplungskontrolle\\
am Wendelstein 7-X\\[1em]
\large Stichpunkte Leitfaden}
\author{}
\date{\today}

\begin{document}
\section{TITELFOLIE + INHALTSVERZEICHNIS}

\begin{itemize}
\item Bevor wir anfangen möchte ich gerne den Unterton der nächsten 90min vorgeben: Zitat französischer, zeitgenössischer Philosophen Guy-Man de Homem-Christo und Thomas Bangalter: "It's been much too long"
\item Willkommen und vielen Dank für Ihr zahlreiches Erscheinen zu meinem Promotionskolloquium.
\item Titel meiner Arbeit: ..., zwar etwas ausladent, aber es soll heute im Besonderen um die Bolometer-Diagnostik am Stellarator Wendelstein 7-X und die von mir eingeführte Echtzeit-Auswertung, sowie tomographische Analyse von synthetischen und experimentellen Daten gehen
\item \textbf{KLICK}
\item Vortrag stelle ich zentrale und kritische Herausforderungen sowie Aufgaben meiner Arbeit vor: Design, Implementierung und Einsatz eines Bolometer-basierten Echtzeit-Feedbacks, die parametrisierte Optimierung und Sensitivitätsanalyse der Kanalauswahl sowie Validierung und Anwendung einer optimierten Tomographie-Methode für das Bolometer-System an W7-X
\end{itemize}

\section{MOTIVATION}

\subsection{Kernfusion: Energieherausforderung}

\begin{itemize}
\item Energiegewinnung aus D-T-Kernfusionen in Hochtemperaturplasmen ist sauber und Rohmaterialien sind weithin ausreichend verfügbar - Deuterium aus Meerwasser, Tritium aus Lithium-Mantel nach Initialzünding
\item Notwendiges Niveau an Effizienz und Plasmanergie durch Lawson-Kriterium bzw. Tripleprodukt definiert: $ n_e\cdot\tau_E\cdot T\geq3\times10^{21} $ keVs/$m^3$ für Netto-Energiegewinn Q>1
\item bei Magnetischer Einschluss $\sim$2,5T bei W7-X erfordert Grössenordnungen T $\sim$ 10-20 keV (100-200 Millionen °C) und Einschlusszeit $\tau_E\sim$ 1-10s
\item \textbf{KLICK}
\item Herausforderung: Management der Wärmelasten auf plasmazugewandten Komponenten während langer Pulse (bis zu 30 Minuten an W7-X geplant)
\item Maschine und Plasma-zugewandte Oberflächen wie besonders Divertor muss einerseits aktiv gekühlt werden - nach OP1.2b - und Plasmar(rand) optimiert werden um Hitzeintrag zu minimieren
\item \textbf{KLICK}
\item Prozess Detachment - Ablösen des Plasma-/rands vom Divertor wobei kinetisch und durch Strahlung ein Grossteil der Energie abgebaut wird
\item \textbf{KLICK}
\item Manipulation der Plasmaprofile, speziell im Rand/SOL um Leistungsverlust (Abstrahlung) und Transport über Separatrix hin zum Divertor zu minimieren
\item Bild: Strahlungseffizient aus Linienstrahlung für verschiedene Verunreinigung über Temperatur
\item Strahlungskühlung $P_{rad}/(V\,n_e\,n_imp)$ hat Maxima bei unterschiedlichen Temperaturen für jedes Element: C Maximum $\sim$10-50\,eV (Rand), O Maximum $\sim$100\,eV $/rightarrow$ Kohlenstoff Potential für Detachment
\item beides intrinsische Verunreinigungen vorranging aus Wandwechselwirkungen - Kohlenstoff dominiert um Faktor $10^2\times$ über Sauerstoff in relevanten Dichte- und Temperaturbereichen im Rand
\item \textbf{KLICK}
\item genereller Ansatz: Verunreinigungseinspeisung - Injektion von bspw. niedrig-Z (He, N$_2$) Elementen zur Steigerung der Strahlung im Plasmarand, Einbringen extrinsischer Verunreinigungen
\item \textbf{KLICK}
\item ohne Akkumulation dieser im Plasmakern -> Störung der Dynamik und pontentieller Strahlungskollaps
\item Erster Einsatz erfordert Validierung und zum Schutz der Maschine und besseres Verständnis: Injektion und Feedback mit Arbeitsgas, so Profiländerung, Kohlenstoffkühlung
\item Erfordert Echtzeitkontrolle, potentiell mit Strahlungsverlust - Motivation für Bolometer-Feedbacksystem
\end{itemize}

\section{ECHTZEIT-FEEDBACK}

\subsection{System}

\begin{itemize}
\item Konzept nicht neu, Durchgeführt an anderen Forschungsreaktoren und geplant für ITER, DEMO
\item \textbf{KLICK}
\item Annahme: Bolometer ist Mess-Diagnostik und Aktuator ist Piezoventil-Box für thermischem Gaseinlass
\item Sehen hier stark reduziertes Modell von W7-X, nur Andeutung der 5 vollsymmetrischen Module inkl. Divertoren oben und unten, unterbrochen jeweils durch Dreiecksebene in der sich das Bolometer befindet
\item Zwei Gas-Boxen inkl. Verunreinigungs-Spektroskopie und Schema der  umlaufenden Feldlinien auf welchen injiziert wird
\item Separatrix und X-Punkte sind Strahlungs-Hotspots: höhere Emissivität aufgrund von Linienstrahlung der Verunreinigung wie vorher erkannt und Plasmaprofilen dort (niedrigen Temperaturen)
\item Sichtlinien der Multikamerasystems des Bolometer in der Dreiecksebene decken X-Punkte, Separatrix und mag. Inseln vollständig ab für optimale Sensitivität gegenüber Ablösungsübergängen
\item \textbf{KLICK}
\item Feedback-System nicht exklusiv für Bolometer: Prozess generalisiert und anwendbar auf mehrere Variablen
\item System verwendet für die detaillierte Steuerung über alle Plasmaparameter einen PID bzw PIDs - Standard in Anlagenkontrolle bzw. in Steuerungstheorie
\item Rechts Piezoventilbox, in Divertornähe Gaseinlass inkl. Schema der Spektroskopie-Sichtlinien mit räumlich-radialer Auflösung
\item \textbf{KLICK}
\item Zwei Vorhersage-Proxies entwickelt: $P_{pred}^{1}$ geometriegewichtet (nutzt mehrere Kanäle mit räumlichen Gewichten), $P_{pred}^{2} $ Einzelkanal (einfachste Implementierung)
\item \textbf{KLICK}
\item Beispiel-Kanalauswahl: VBC S={61,65,73,78,86} erreicht zBsp. ohne a priori Wissen $\geq$85\% Vorhersagegenauigkeit, beobachtet Separatrix/SOL-Region im speziellen
\item \textbf{KLICK}
\item Zielgrösse der Vorhersage nicht direkt $P_{pred}$ oder $P_{rad}$ sondern $f_{rad}$ als eigentliche Kontrolgrösse
\end{itemize}

\subsection{XP20181010.32 Erfolge}

\begin{itemize}
\item Prozess Verlauf natürlich in meheren Schritten - mit mehr und weniger Erfolg
\item Durchbruch-Experiment: stabile 'Hochleistungs'-Feedback-Entladung mit Wasserstoff-Einlass
\item Dauer: $\sim$9s Feedback-Kontrolle, $P_{ECRH}$=6.25\,MW, Gesamtenergie 55\,MJ
\item \textbf{KLICK}
\item Zielpunkt Strahlungsverlust mit $P_{pred}^{1}$ $f_{rad}\sim$90\% aufrechterhalten, Spitzen bis 100\% während übergängen
\item \textbf{KLICK}
\item Vorhersagegenauigkeit: $P_{pred}^{1}$ verfolgt $P_{rad}$ innerhalb 10-15\% über gesamte Entladung, validiert Proxy-Konzept
\item Sogar $P_{pred}^{2} $ mit ähnlicher Genaugikeit zu $P_{rad}$ ohne Dimensionierung und Ableitung
\item \textbf{KLICK}
\item Plasmanergie: Unter konstanter Heizleistung, $W_{dia}\sim$ stabil inkl. minimaler Oszillationen trotz hoher Strahlung aufrechterhalten, gute Einschlusszeit angedeutet
\item Vorhersage-$f_{rad}$ überschätzt 'tatsächlichen' Verlustwert konstant $\sim$10\%
\item Divertor Wärmelast-Reduktion: Faktor $\geq2\times$ erreicht (von $\sim$4\,MW auf <2\,MW)
\item \textbf{KLICK}
\item Strahlungsfeedback kann $f_{rad}\sim0.9-1$ inkl. Oszillationen aufrechterhalten und Divertor-Wärmelast signifikant reproduziert
\item \textbf{KLICK}
\item Spektroskopie-Messung exklusiv für C$^2$ Linienstrahlung zeigt Ablösung; Erscheint ab $f_{rad}\sim$50\%, vollständige Ablösung bei $f_{rad}$>80\% mit Oszillationen hin zu 100\% und weiterer Verschiebung
\item \textbf{KLICK}
\item Zusammenfassung: Demonstriert Szenario sehr gut anwendbar, Langzeitpotential, hohe Strahlung ohne Einschlussdegradation
\end{itemize}

\section{LOS-SENSITIVITÄT}

\subsection{LOS-Sensitivität (Motivation)}

\begin{itemize}
\item Wg. Experimentplan und anderer Anforderungen an das System keine Zeit für physikalische Untersuchung und weitere Modellierung für $P_{(pre)}$
\item Anschliessende Frage: Welche Bolometer-Kanäle sind am sensitivsten für Plasmaparameteränderungen oder am besten geeignet für Feedback?
\item Ziel: Optimale Kanalauswahl für Echtzeit-Feedback-Kontrolle
\item \textbf{KLICK}
\item Systematischer Ansatz: Räumliche Kartierung der Sensitivität,
\item \textbf{KLICK}
\item Datensatz: 45+ Entladungspaare aus OP1.2b, Feedback Experimente - nicht nur $P_{rad}$
\item \textbf{KLICK}
\item 5 Evaluierungsmetriken entwickelt, um verschiedene Aspekte der Kanalsensitivität zu quantifizieren + Variation der Grösse der Sichtlinien-Auswahl
\item Analyse-Workflow: Für jede Entladung, berechne Metrik für jede mögliche Kanalkombination und Kombinationsgrösse, aggregiere über Kombinationen und Integriere für zentrales Sensitivitätsprofil
\item \textbf{KLICK}
\item Bilder beschreiben: Sensitivität vs. $\rho_{\text{pol}}$ ...
\item \textbf{KLICK}
\item Allgemein: robuste Kanalauswahl jederzeit $\sim$80\% Vorhersagegenauigkeit mit kleinen Fehlerbalken (m=5)
\item \textbf{KLICK}
\item Aussage über eine einzelne, hervorverstechende Kanalauswahl nicht möglich oder existiert nicht in diesem Rahmen
\item \textbf{KLICK}
\item Kernkanäle ( $\rho_{\text{pol}}$<0,8) zeigen niedrigere Sensitivität aufgrund flacher Strahlungsprofile - weniger Information über Plasmaänderungen
\item \textbf{KLICK}
\item 5 oder mehr Kanäle hinreichend für akkurates $P_{pred}$ bzw. Feedback-Zwecke
\item \textbf{KLICK}
\item Optimale Kanalanzahl: Systematische Variation m=3-7 zeigt m=3-5 ausreichend für Feedback (Kompromiss zwischen Redundanz und Komplexität)
\item : für HBC Sichtlinien über Separatrix/SOL ($\rho_{\text{pol}}\sim$\,0,95-1,05), bzw. für VBC auf X-Punkte konzentriert zeigen höchste Sensitivität über alle Metriken

\end{itemize}

\section{TOMOGRAPHIE}

\subsection{Motivation}

\begin{itemize}
\item Bisher: 1D linienintegrierte Messungen gezeigt und verrechnet
\item Multikamerasystem - Inversion um Verlauf der gesuchten Emissivitätsverläufe besser zu verstehen
\item Ziel: 2D-Strahlungsverteilung $P_{rad}$ (R,Z) aus Strahlungsmessungen rekonstruieren
\item \textbf{KLICK}
\item Herausforderung: Invertierung ist unterbestimmtes, lineares System - 61 Messungen oder (24x2)x32 LOS-Kreuzungen, $\sim$4500 Gitterzellen
\item T allgemein nicht quadratisch, keine Inverse und schlechte bzw. hohe Kondition
\item \textbf{KLICK}
\item Lösung: Regularisierungs-Ansatz durch Fisher Information
\item Diskretisierbares und numerisch lösbares Minimal-Problem mit zusätzlichen Stellschrauben, besonders in $\mathbf{H}$
\item Fischer-Info misst Erwartungswert des Informationsgehaltes in $\vec{x}$ in $\vec{r}$
\item Regularisierungsoperator links-rechts Laplace
\item Zentrale Beziehung zur Cramer Rao Ungleichung: deckelt die Varianz der Informations-Ungenauigkeit durch Regularisierung
\item MFR-Prinzip: Maximiere Entropie + Glätte bei Erfüllung von Randbedigungen - bevorzugt physikalisch plausible (glatt entland Flussflächen) Lösungen
\end{itemize}

\subsection{MFR mit RDA}

\begin{itemize}
\item Gehen noch einen Schritt weiter: welche zusätzlichen a prior Informationen können wir vor dem Hintergrund der Feedback Experimente einfließen lassen
\item Radiale Anisotropie-Gewichtung: $k_{ani}(\rho)$ zusätzlicher Faktor auf poloidale Ableitung (<1 individiülle Beiträge gewichtiger, >1 weniger Einfluss) variiert radial
\item Physikalische Motivation: Strahlung folgt Flussflächen im Kern (isotrop), wird (potentiell) anisotrop(er) im SOL/Divertor kürzere Verbindungslängen, offen Flussflächen, Inseln etc.
\item \textbf{KLICK}
\item Gleichungsmodifikation
\item \textbf{KLICK}
\item Beispielprofil
\item \textbf{KLICK}
\item Allgemein glatte Resultate + störungsresistent, numerisch stabil und parametrisierbare An-/Isotropien

\item \textbf{KLICK}
\item Zusammenfassung der Störungrechnung an Grundlage des Kamerasystems
\item \textbf{KLICK}
\item Diskretisierung der Sichtlinien bzw. Detektoren und der Tomographie-Domäne ist allgemein numerisch stabil
\item \textbf{KLICK}
\item Vorwärtsrechnung mit vollsymmetrischen Verteilung ergab kleine aber nicht vernachlässigbare Asymmetrie in linienintegrierten Messungen
\item \textbf{KLICK}
\item Nachträglich gemessene Störung an Design-Geometrie ergibt Fehler mit selber Größenordnung wie eben diese, aber keine Aussage über wann/wie/wo
\item \textbf{KLICK}
\item Gerichtete Variation ergibt worst case
\end{itemize}

\subsubsection*{Phantoms}

\begin{itemize}
\item Phantom-Benchmarks: Regularisierungen müssen/sollten mit a priori Wissen untersucht und folglich eingestellt werden um sinnvolle bzw. verlässliche Ergebnisse zu erhalten
\item \textbf{KLICK}
\item synthetische Tests mit verschiedenen Strahlungsmustern (Kern-dominiert, Rand-dominiert, gemischt)
\item Metriken um Qualität zu messen: Rekonstruktionsqualität $\rho_c$ (Korrelation mit Wahrheit), $\chi^2$ (Anpassungsgüte), zweidimensionale, mittlere Abweichug (MSD), 2D integrierte Leistungen
\item \textbf{KLICK}
\item nur $\chi^2$ für experimentelle Rekonstruktionen verfügbar; außerdem nur für eine magnetische Konfiguration und Design-Geometrie
\item \textbf{KLICK}
\item Beispiel: helle, Insel-artige Emissivität in SOL mit weniger intensivem, isotropem Ring in Kern
\item \textbf{KLICK}
\item Experimentel motivierte Gewichtung: scharfe Anisotropie und schwacher Ring innen
\item \textbf{KLICK}
\item Ratio hin zu glatten Strukturen, maximale Helligkeit reduziert und Anisotropien entlang poloidaler Schalen des Gitters vergrößert
\item \textbf{KLICK}
\item Skalierung um integralen Faktor führt zu kristalisierung an den vorher festgestellten Punkten; $k_{ani}\gg\Delta r/\varphi$ erzwingt Festsetzung an $max(I_{F})$
\item \textbf{KLICK}
\item Statistik: Beispielhafte Ergebnisse für nur $k_{edge}$ Variation - Klares Maximum in Pearson nicht kongruent mit $\chi$ und MSD!
\item Legt ein grundlegendes Problem nahe: Maß an $\chi^2$ nicht zwangsweise Beste Metrik für Optimierung
\item diese Optimierung evtl. nur für Klassen von Verteilungen/Regularisierungsparametern gültig und nötig, nicht allgemein für alle Emissivitäten valide Ergebnisse
\item \textbf{KLICK}
\item Integrale Statistik über all angewandten synthetischen Profile
\item Qualitätsparameter: Pearson\,$\ne\chi^{2}$ aber stimmt mit MSD Verlauf überein
\item $\chi^{2}$ als Optimierungsgrösse der Tomographie genügend aber durch Faltung mit Fehlerbalken und Geometrie-Einflüssen nuanciert zu betrachten
\item Zusammenfassung
\item Zusätzlich: Vergleich Virtueller Kameramessungen mit 2D Strahlungsintergralen zeigt Diskrepanz die in etwa Größenordnung SOL Emissivität ist
\item \textbf{KLICK}
\item Diskrepanz Qualitätsparameter
\item \textbf{KLICK}
\item Optimierung $k_{ani}$ hochdimensional
\item \textbf{KLICK}
\item Ansatz validiert und robuste, wertvolle Ergebnisse erzielt
\item \textbf{KLICK}
\item Folgende experimentelle Rekonstruktion fußen auf die gemachten Erkenntnisse
\end{itemize}

\subsubsection*{Experimentelle Rekonstruktionen}

\begin{itemize}
  \item Zur Veranschaulichung der gewonnen Erfahrung ein praktisches Beispiel
  \item LBO Events: Verunreinigungs-Injektion durch Laser-Ablationen von Metall auf Fenster ins Plasma; hohe Temperatur und kinetische Energie
  \item Ionisations-Kaskaden inkl. Transport in eingeschlossenen Bereich erwartet
  \item Starke Strahlungsspitzen, sonst konstante Plasmaparameter
  \item Fokus auf ersten Laser-Schuss bei 2s, Strahlungsantwort verzögert
  \item \textbf{KLICK}
  \item direkt: neues Gitter, andere magnetische Konfiguration (auch Teil der vorheringen Iteration und Rechnungen)
  \item Mit Erfahrung aus Phantomen qualitativ und quantitativ erfolgreiche, verlässliche Rekonstruktionen, Strahlung stark im inneren, untersten X-Punkt und entlang der Separatrix
  \item \textbf{KLICK}
  \item angesprochene Statistik der Leistungen: hier als einzige Parameter Verfügung vgl. Phantome
  \item Diskrepanz P$_{rad}$ vs P$_{2D}$ reproduziert mit SOL/Kern Korrelationen bis auf die genauen Werte
  \item Bestätigung mit Phantomen und synthetischen Daten unterstützt die Ergebnisse und validiert die 2D Leistung aus Tomogramen als (mehr) angemessenen Strahlungsverlustwert
\end{itemize}

\subsubsection*{Globale Leistungsbilanz}

\begin{itemize}
\item Letzter Ausblick der Anwendungsmöglichkeiten mit sekundären Ergebnissen aus Tomographie
\item Simpler 0D Ansatz aus quasi Equlibrium für globale Leistungsbilanz aus gespeicherter Energie, Eintrag und Verlusten;
\item \textbf{KLICK}
\item XP20181010.32 als bekannte Referenz nochmal heranziehen
\item \textbf{KLICK}
\item Vergleich P$_{bal}$ aus P$_{rad}$ und getrennten P$_{2D}$ Gesamt und nur aus Plasma-Kern
\item Bessere Repräsentation und stabiler bei Plasma-übergängen, allgemein näher an P$_{bal}=0$ für Emissivität nur aus Kern
\item \textbf{KLICK}
\item Zusammenfassung
\item Korrelation aber nicht Kongruenz der Kamera- und Verteilungs-Strahlungsleistungen
\item \textbf{KLICK}
\item Trends und qualitative Aussagen aber in Übereinstimmung mit Ergebnissen aus surrogaten Daten
\item \textbf{KLICK}
\item 2D Leistung generell unter- und SOL Helligkeit überschätzt
\item \textbf{KLICK}
\item Schlussfolgernd: Hinreichende Ergenisse mit globaler Leistungsbilanz
\end{itemize}

\section{SCHLUSSFOLGERUNGEN}

\subsection{Haupterfolge}

\begin{itemize}
\item Drei Hauptbeiträge
\item \textbf{KLICK}
\item (1) Echtzeit-Feedback: Erste Bolometer-basierte Implementierung am W7-X, erreichte $f_{rad}\geq$ 85\% mit Faktor $\geq2\times$ Divertor-Wärmelast-Reduktion
\item \textbf{KLICK}
\item (2) LOS-Optimierung: Systematische Analyse von 45+ Entladungspaaren identifizierte optimale Kanalauswahl, m=5 Kanäle ausreichend für robuste Kontrolle, besonders Randkanäle bzw. auf Inseln/X-Punkte
\item \textbf{KLICK}
\item (3) Tomographie: MFR+RDA-Methode validiert durch 12+ Phantom-Tests und experimentelle Entladungen, liefert 2D-Strahlungsverteilung
\end{itemize}

\section{ABSCHLUSS}

\subsection{Danksagung}

\begin{itemize}
\item Vielen Dank für Ihre Aufmerksamkeit - "Sin é, a chairde" (Gälisch: "That's it, folks.")
\item Fragen, bitte immer her damit!
\item Danksagung: W7-X-Betriebsteam, IPP-Greifswald-Kollegen, Prof. Dr. Klinger, der Arbeitsgruppe 'Transport, Verunreinigungen und Strahlung' um Dr. Felix Reimold und Dr. Rainer Burhenn und Familie
\end{itemize}


\pagebreak%
\section{AUSBLICK}

\begin{itemize}
\item \textbf{Upgrades:} (1) Latenz für schnellere Feedback-Antwort verringern (2) Prädiktive Algorithmen zur Antizipation von Strahlungsvorhersagen, maschinelles Lernen für adaptive Kontrolle,  automatische Kanalauswahl basierend auf Plasmazustand (3) Konfigurationsraum ausnutzen mit Skalierungszahlen (4) Gas-Injektions-Timing-Optimierung
\item \textbf{KLICK}
\item \textbf{Modellierung:} (1) bspw. EMC3-EIRENE für Erforschung der Feedback Szenarios nutzen und Phantom-Untersuchung zu erweitern (2) vollständige Parametrisierung der Phantom-Ergebnisse und Sichtlinien-Optimierung
\item \textbf{KLICK}
\item \textbf{Tomographie-Erweiterung:} (1) Erweiterte Phantom-Bibliothek (>50 Tests), (4) GPU-Beschleunigung für <100ms/Frame-Rekonstruktion (5)Bayes'sche Optimierung oder Gittersuche für $ k_{an}i ( \rho $ )-Profil, adaptive Anisotropie basierend auf Plasmazustand
\end{itemize}

\subsection{Validierung}

\begin{itemize}
\item Alle vorherigen und folgenden Untersuchungen basieren auf Design des Bolometers (idealisierte Geometrie per Bauplan ohne 'Fehler')
\item Geometrische Untersuchungen: Diskretisierung des Bolometers ohne signifikanten Einfluss
\item Beihilfe mit synthetischen Tests aus Phantom-Strahlungsprofilen aus bspw. STRAHL mit experimentellen Eingangsdaten
\item Für verschiedene Strahlungsniveaus aus dem Feedbackexperiment XP20181010.32
\item \textbf{KLICK}
\item Voll symmetrische Emissivität und Design-Geometrie (HBC) ergeben leichte Asymmetrie in linienintegrierten Messungen
\item Fragestellung: Welche Kamerageometrie ist denn symmetrisch bzw. woher kommt die Asymmetrie?
\item Künstliche Geometrie erarbeitet: Zentrierung des Arrays, Spiegelung der Sichtlinien innerhalb der Kamera und schliessliche 'Aufrechtstellung' und Ausrichtung genau in die 108$^\circ$ Ebene
\item \textbf{KLICK}
\item Vergleichbar fast vernachlässigbare änderung in $\mathbf{T}$
\item \textbf{KLICK}
\item Aber voll symmetrisches Ergebnis in $P_{M}$
\item \textbf{KLICK}
\item Post-OP1.2b Roboter-gestützte Messungen der HBC Kameraöffnung zeigen Verschiebungen, potentiell durch thermischen und mechanischen Stress auf die Maschine
\item Unbekannt wie der Zustand während der Messkampagne ist: Anwendung mit Störungsrechnung dieser Abweichungen auf Kamerageometrie zeigt bis zu $\pm$2$^{\circ}$ Abweichung in Kamera-Ebene als Worst-Case-Szenario
\item \textbf{KLICK}
\item Integration des vorherigen Phantoms mit solcher Geometrie gibt $\pm$0.2r$_{a}$ Verschiebung im Vergleich zu symmetrischen Ergebnissen
\item \textbf{KLICK}
\item Hypothetischer Fall: wissen nicht von der Verschiebung bzw. Verzerrung der Geometrie - Messung mit gestörter Kamera-Ausrichtung und Tomographie mit Design-Geometrie (Blindfehler)
\item \textbf{KLICK}
\item Zusammenfassung
\end{itemize}


\subsection{Zusammenfassung}

\begin{itemize}
\item \textbf{Feedback-System:} 13,6ms Akquisitionslatenz (minimal), 150-400ms Gesamtschleife, PID-Regler mit dual-Proxy-Vorhersage
\item \textbf{Benchmark-Experiment:} XP20181010.32 - 9,2s Daür, $P_{ECRH}$ =6,23 MW, $f_{rad} \geq $ 85\% aufrechterhalten, Faktor $\geq2\times$ Wärmelast-Reduktion, keine Disruption
\item \textbf{Diagnostik-Leistung:} 94 Kanäle installiert (61 funktional in OP1.2b), Kalibrierungsstabilität <5\% über 1182 Experimente, SNR>1000
\item \textbf{LOS-Sensitivität:} Separatrix/SOL-beobachtende Kanäle ( $\rho_{\text{pol}} \sim$ 0,95-1,05) optimal, VBC übertrifft HBC für Randsensitivität, mehrere Kanal-Sets erreichen $\geq$ 85\% Genauigkeit
\item \textbf{Ablösungsphysik:} C$^{2+}$ -Signatur erscheint bei $f_{rad}\sim$50\%, vollständige Ablösung bei $ f_{rad} $ >90\%, kritische Region bei $\rho_{\text{pol}}\approx$ 0,95
\item \textbf{STRAHL-Modellierung:} Kohlenstoff dominiert um Faktor 10$^2\times$ über Saürstoff, C$^{3+}$/C$^{2+}$ -Verhältnis kritisch an LCFS für Strahlungsverteilung
\item \textbf{Tomographie-Qualität:} Rekonstruktionen erreichen $\rho_c$ >0,85, $\chi^2$ <2,5, $ P_{2D} $ innerhalb 10-20\% von $P_{rad}$ , zeigen X-Punkt-Lokalisierung und Inselketten-Struktur
\item \textbf{Leistungsbilanz:} $P_{2D}^{Kern} $ stabil für quasi-stationäre Phasen, Standard-MFR unterschätzt um 15-25\% (SOL-Beitrag fehlt)
\item \textbf{Validierung:} OP1.2b-Kampagne 2018, 1182 Experimente analysiert, 45 Entladungspaare für LOS-Studie, 4 für Tomographie-Validierung

\item \textbf{Andere Geräte:} Anwendbar auf LHD (Stellarator), ITER, DEMO, SPARC (Tokamaks) - universelle Strahlungskontroll-Strategie
\item \textbf{Vision:} Vollautomatisches, adaptives, prädiktives Strahlungsmanagement-System für zukünftige Fusionsreaktoren
\item \textbf{ITER/DEMO-Anwendung:} Direkt übertragbare Methoden, $P_{fus}\sim$ 500 MW, $f_{rad}\geq$95\% erforderlich, $q_{di}$<10 MW/m$^2$ Ziel
\item \textbf{Kammermodelle:} Verfeinerung von Zwei/Drei-Kammer-Modellen für Verunreinigungstransport, Entwicklung von Skalierungsgesetzen für verschiedene Betriebsregime
\end{itemize}

\section{BACKUP-FOLIEN - Übersicht}


\begin{itemize}
\item Backup-Folien liefern detaillierte technische Information für tiefgehende Fragen
\item Abgedeckte Themen: Plasmaphysik, STRAHL-Modellierung, Diagnostik-Details, LOS-Sensitivitätsmetriken, Tomographie-Vergleiche, DAQ-Latenz
\item Verfügbar, um spezifische Fragen zu Methodik, Implementierungsdetails oder alternativen Ansätzen zu adressieren
\end{itemize}

\vspace{1em}


\section{MOTIVATION}

\subsection{Kernfusion: Energieherausforderung}

\begin{itemize}
\color{ipp}{
  \item Transport am W7-X: neoklassische Basislinie mit anomalen Beiträgen von Turbulenz (ITG, TEM-Moden)
  \item Anomaler Transport erhöht Diffusion um Faktor 10-100 über neoklassisch, kritisch für Verunreinigungstransport-Modellierung
  \item Strahlungsfraktion $ f_{rad} = P_{rad} / P_{ECRH} $ muss ausbalanciert sein: $ f_{rad} $ <70\% unzureichend für Divertorschutz, $ f_{rad} $ >95\% riskiert Kernkühlung und Einschlussdegradation
  \item Optimales Betriebsfenster: $ f_{rad} $ =85-90\% erreicht Ablösung bei Erhalt der Plasmastabilität
}
\end{itemize}

\section{KERN-BOLOMETER AM W7-X}

\subsection{Bolometer-Diagnostik (Hardware)}

\begin{itemize}
\item Metallwiderstand-Design: dünner Goldfilm (1,5 $\mu m $ ) auf Glimmersubstrat, absorbiert Strahlung $\rightarrow $ Temperaturerhöhung $\rightarrow $ Widerstandsänderung
\item Wheatstone-Brücken-Konfiguration: AC-Anregung (5V, kHz) eliminiert 1/f-Rauschen, differentielle Messung verdoppelt Signal
\item Drei Kamera-Arrays bieten umfassende Plasmaabdeckung: HBC (32 Kanäle, horizontal), VBCl/r (2 $\times $ 31 Kanäle, vertikal)
\item Räumliche Auflösung $\sim $ 5 cm an magnetischer Achse, bestimmt durch Etendü-Breite und Beobachtungsgeometrie

\end{itemize}

\subsection{Leistung}

\begin{itemize}
\item Bolometer messen totale Plasmastrahlungsleistung (0,2-600 nm) durch resistive Heizung von Metallabsorbern
\item Absorptionseffizienz nahe 1 (>95\%) über 0,2-600 nm Wellenlängenbereich - breitbandige Messung
\item Zeitauflösung konfigurierbar: {0,4, 0,8, 1,6, 3,2, 6,4, 12,8} ms Abtastung verfügbar, 1,6ms für Feedback verwendet (balanciert Geschwindigkeit vs Rauschen)
\item Kalibrierung ist kritisch für absolute Leistungsmessungen - drei Methoden verwendet: (1) Laser, (2) elektrische ohmsche Heizung, (3) in-situ während Betrieb
\item Leistungsmetriken: SNR>1000 bei hoher Strahlung, Rauschen $\sigma_{\Delta U}=0,339 \pm 12,586 \mu V $ , Gauß-Fehler $ 1,6 \mu W $ , Drift $ 0,057 \pm 27,621 \mu $ V/s
\item Jeder Kanal integriert Strahlung entlang der Sichtlinie - nutzt Geometriematrix T und Etendü $ K_M $ zur Rekonstruktion der 2D-Verteilung
\item Sehnen-Helligkeit $ P_{ch} = P_M/K_M $ liefert radiale Profilinformation - zeigt Strahlungsverteilung von Kern bis SOL
\item Etendü $ K_M $ quantifiziert geometrische Lichtsammeleffizienz: berücksichtigt Detektor-Apertur-Geometrie, partielle Abschattung, winkelabhängige Transmission
\item Bolometer zentrale Plasmadiagnostik (integrale Strahlungsmessung): Einsatz für Leistungsbilanz, tomographische Rekonstruktion der 2D Strahlungsverteilung und nun Echtzeit-Feedback

\color{ipp}{
  \item Kalibrierungsparameter (OP1.2b Median): $R_M=0,987\pm0,011\,$k$\Omega$, $\kappa_M=0,724\pm0,055\,A^2$, $\tau_M=110,57 \pm$4,93\,ms, bemerkenswert stabil über 1182 Experimente
}
\end{itemize}

\section{ECHTZEIT-FEEDBACK}

\subsection{System}

\begin{itemize}
\color{ipp}{
  \item Akquisitionslatenz: 13,6ms minimal (1,6ms Abtastung + 12ms Verarbeitung), gemessen durch Laserdioden-Validierung
  \item FIFO-Glättung: $\sim$ 8ms zusätzliche Verzögerung (M=10 Samples $\times $ 1,6ms Abtastperiode) reduziert Rauschen
  \item Gesamtschleifenlatenz: 150-400ms (plasmadominiert durch Gasventil-Antwort $\sim$50-200ms und Verunreinigungstransport $\sim$100-300ms)
  \item Abtastrate-Kompromiss: 1,6ms optimal balanciert Geschwindigkeit (schnelle Antwort) vs Rauschen (höhere Abtastraten verstärken Rauschen)
  \item Erste Bolometer-basierte Feedback-Implementierung am W7-X (OP1.2b-Kampagne, 2018) - vorherige Systeme nutzten Interferometrie
  \item Hardware: NI 6321 DAQ-Karte (16-bit, 250 kS/s), PID-Regler in LabVIEW implementiert, parallele Verarbeitung für niedrige Latenz
  \item Sollwert: $f_{rad}^{target}=85-90\%$ balanciert Divertorschutz (erfordert $\geq $ 85\%) gegen Kernkühlung-Risiko (vermeidet >95\%)
  \item Laserdioden-Validierung: 13,6ms Akquisitionslatenz experimentell bestätigt durch kontrollierte Lichtpulse
  \item Zukünftige Verbesserungen: <10ms Akquisition möglich mit Hardware-Upgrade, prädiktive Algorithmen könnten Gesamtlatenz auf $\sim$50-100ms reduzieren
}
\end{itemize}

\subsection{XP20181010.32}
\subsection*{Diskussion}

\begin{itemize}
\item \textbf{Bolometer:} Direkte $P_{rad}$-Messung (keine Modellannahmen), breitbandig, schnell (1,6ms Zeitauflösung)
\item \textbf{Robustheit:} Unabhängig von Verunreinigung (misst totale Strahlung), funktioniert mit C, N$_2$, Ne, Ar ...
\item \textbf{Räumliche Information:} Kanäle unterscheiden Kern vs Rand vs SOL-Strahlung, ermöglicht gezielte Kontrolle
\item \textbf{Latenz-Limitierung:} 150-400ms Gesamtschleife begrenzt Bandbreite, anscheinend hinreichend für Ablösungskontrolle aber vorliegende Konfiguration führt zu Oszillationen in Plasma-Antwort-Feedback-Kreislauf
\item \textbf{Breite Anwendbarkeit:} Konzept direkt anwendbar auf ITER, DEMO, und zukünftige Stellaratoren - universelle Lösung
\end{itemize}

\section{LOS-SENSITIVITÄT}

\subsection{LOS-Sensitivität (Motivation)}

\begin{itemize}
\color{ipp}{
  \item Beste Kanäle identifiziert: VBC S={61,65,73,78,86} (beobachten Separatrix/SOL-Region), HBC S={4,7,8,20,23,27} (Rand-Fokus)
  \item Sensitivität korreliert stark mit Strahlungsgradienten - Kanäle, die Regionen mit steilen $\partial P_{rad} / \partial \rho $ beobachten, sind am sensitivsten
  \item Validierung: Experimentelle Feedback-Leistung mit ausgewählten Kanälen bestätigt Vorhersagen - stabile $ f_{rad} $ -Kontrolle erreicht
}
\end{itemize}

\section{PLASMAPHYSIK}

\subsection{Plasmatransport am W7-X}

\begin{itemize}
\item Drei Transportmechanismen: klassisch, neoklassisch, anomal
\item W7-X-Optimierung für niedrige Kollisionalitätsregime
\item Anomaler Transport dominiert um Größenordnung
\item Verunreinigungstransport ladungsabhängig
\end{itemize}

\subsection{Verunreinigungstransport und Strahlung}

\begin{itemize}
\item Anomale Diffusion $ D_{an}\sim $ 0,3-3 m $ ^2 $ /s dominiert
\item Neoklassische Konvektion erzeugt Einwärts-Pinch für hoch-Z
\item DEMO-Ziel: $ f_{rad} $ >0,95 mit 70/30 Kern/SOL-Aufteilung
\item Kontrollherausforderung: Balance Schutz vs Einschluss
\end{itemize}

\subsection{Plasmaverunreinigungen}

\begin{itemize}
\item Kohlenstoff dominant: 10 $ ^2 \times $ höher als Saürstoff
\item Ionisationsstufen variieren mit Temperatur
\item Strahlungsverteilung entwickelt sich mit $ f_{rad} $
\item Kern/SOL-Umkehr bei hoher $ f_{rad} $
\end{itemize}

\subsection{Verunreinigungseffekte und Transportsensitivität}

\begin{itemize}
\item Brennstoffverdünnung beeinflusst Lawson-Kriterium
\item Diffusionssensitivität: 120\% Spitzenanstieg mit niedrigerem D
\item Separatrix-Profil kritisch für Ablösung
\item STRAHL validiert experimentelle Trends
\end{itemize}

\subsection{Plasmaablösung}

\begin{itemize}
\item Schwelle: $ T_e $ < 5 eV am Target
\item Volumetrische Verluste dominieren
\item Strahlungsverstärkungs-Rückkopplungsschleife
\item W7-X erreichte $ f_{rad} $ =100\% Ablösung
\end{itemize}

\section{MODELLIERUNG \& ANALYSE}

\subsection{Verunreinigungseinspeisung-Modelle}

\begin{itemize}
\item Zwei-Kammer- und Drei-Kammer-Modelle
\item Erfassen grundlegende Zeitskalen: Wand, Plasma, SOL
\item R $ ^2 $ >0,85 Korrelation mit Experimenten
\item Zukunft: Integration mit 3D-Codes
\end{itemize}

\subsection{Plasma-Leistungsbilanz}

\begin{itemize}
\item Energieerhaltung: $ P_bal\approx $ 0
\item Kombinierte Unsicherheiten $\pm $ 1-2 MW
\item Gefäßwärmelast 15-25\% von $ P_{rad} $
\item Tomographie verbessert Stabilität
\end{itemize}

\section{DIAGNOSTIK-DETAILS}

\subsection{Detektorsensitivität: Etendü}

\begin{itemize}
\item $ K_M $ quantifiziert Lichtsammeleffizienz
\item Einheiten: mm $ ^{-3} $ (inverse Volumen)
\item Räumliche Auflösung $\sim $ 4-5 cm an Achse
\item Kritisch für absolute Leistungsmessungen
\end{itemize}

\subsection{Temperaturdrift-Einfluss}

\begin{itemize}
\item Hauptsystematischer Fehler in langen Pulsen
\item Graukörperstrahlung skaliert mit T $ ^4 $
\item Lineare Korrektur für langsame Drifts
\item Zukunft: temperaturabhängige Kalibrierung
\end{itemize}

\subsection{In-Situ-Kalibrierung}

\begin{itemize}
\item Essentiell wegen $\pm $ 5-10\% Variationen
\item Drei Parameter: $ R_M , \tau_M , \kappa_M $
\item OP1.2b: außergewöhnliche Stabilität über 1182 Experimente
\item Kabelkorrektur verhindert systematische Fehler
\end{itemize}

\subsection{Detektorleistungs-Statistiken}

\begin{itemize}
\item SNR>1000 bei hoher Strahlung
\item Rauschen: 0,339 $\pm 12,586 \mu $ V
\item Minimale Degradation beobachtet
\item Erfüllt alle fusionsgerechten Anforderungen
\end{itemize}

\subsection{Bolometer-Gleichungs-Ableitung}

\begin{itemize}
\item AC-Anregung entfernt 1/f-Rauschen
\item Wheatstone-Brücke: $\Delta U \propto \Delta R_M $
\item Leistungsbilanz: Heizung vs Kühlung
\item Trade-off: $\tau und \kappa $ invers verwandt
\end{itemize}

\subsection{Messalgorithmus}

\begin{itemize}
\item Drei Stufen: Initialisierung, Kalibrierung, Messung
\item Echtzeit-Feedback via NI 6321 DAQ
\item FIFO-Puffer verhindert Datenverlust
\item Latenzoptimierung: 13,6ms Minimum
\end{itemize}

\section{STRAHL-MODELLIERUNG}

\subsection{STRAHL-Grundlagen}

\begin{itemize}
\item 1D-Verunreinigungstransport-Code
\item Flussflächen-gemittelte Größen
\item Zeitabhängige Evolution
\item Löst Kontinuität für jedes Z
\end{itemize}

\subsection{STRAHL-Verunreinigungstransport-Modellierung}

\begin{itemize}
\item Leitende Gleichung mit D und v
\item Kohlenstoffdominanz validiert
\item Kritische Ladungszustände: C $ ^{3+} $ /C $ ^{2+} $
\item Strahlungsverschiebung mit $ f_{rad} $
\end{itemize}

\subsection{STRAHL-Parametervariationen}

\begin{itemize}
\item Systematische Sensitivitätsstudie
\item Kernstrahlung stabil $\pm $ 10\%
\item Rand variiert $\pm $ 30-50\%
\item C $ ^{3+} $ /C $ ^{2+} $ am empfindlichsten
\end{itemize}

\subsection{STRAHL: Strahlungsfraktions-Abhängigkeit}

\begin{itemize}
\item $ f_{rad} $ =90\% vs 100\% Vergleich
\item Spitze verschiebt sich von SOL zu LCFS
\item Kern/SOL-Verhältnis ändert sich 0,3 $\rightarrow $ 0,7
\item Validiert Feedback-Kanalauswahl
\end{itemize}

\section{LOS-SENSITIVITÄTS-DETAILS}

\subsection{Kamerageometrie-Sensitivität}

\begin{itemize}
\item Aperturpositionierungs-Unsicherheiten
\item Vorwärtsmodellierung <1\% Unterschied
\item Fertigungstoleranzen akzeptabel
\item Validiert Ingenieurdesign
\end{itemize}

\subsection{Segmentierung und Sichtlinien-Kegel}

\begin{itemize}
\item N=8 optimal für Genauigkeit vs Berechnung
\item Voller Kegel vs infinitesimale Projektion
\item Richtige Segmentierung essentiell
\item Jedes Segment gewichtet durch Etendü
\end{itemize}

\subsection{Tomographie: RDA vs RGS Vergleich}

\begin{itemize}
\item RDA: absolute Differenzen, erhält Gradienten
\item RGS: relative Gradienten, verstärkt Merkmale
\item RDA gewählt für Robustheit
\item RGS nützlich für Post-Shot-Analyse
\end{itemize}

\subsection{Tomographie: Künstliche Kameras}

\begin{itemize}
\item VBCm und MIRh konzeptülle Designs
\item 4-Kamera verbessert um 15-20\%
\item Trade-off: Kosten vs Verbesserung
\item Aktülle 3-Kamera ausreichend
\end{itemize}

\subsection{Phantom-Tomographie: Grenzfälle}

\begin{itemize}
\item Invertierte Anisotropie testet Grenzen
\item Homogene Verteilung validiert keine Artefakte
\item Grenzfälle offenbaren Beschränkungen
\item Leiten Betriebsparameter
\end{itemize}

\subsection{LOS-Sensitivitätsbewertung (Kreuzkorrelation)}

\begin{itemize}
\item Zeitliche Ähnlichkeitsmetrik
\item Große Fehlerbalken: $\pm $ 40\%
\item X-Punkt-Kanäle am schlechtesten
\item Kritisch für Echtzeitkontrolle
\end{itemize}

\subsection{LOS-Sensitivitätsbewertung (Mittlere Abweichung)}

\begin{itemize}
\item Varianz-basierte Qualität
\item Fehlerbalken 10 $\times $ größer als gewichtet
\item Kernkanäle konsistent
\item Rand optimal für Feedback
\end{itemize}

\subsection{LOS-Sensitivitätsbewertung (FFT-Korrelation)}

\begin{itemize}
\item Spektralanalyse-Ansatz
\item Identifiziert oszillatorisches Verhalten
\item Ergänzt Zeitbereich
\item Leitet Reglerentwurf
\end{itemize}

\subsection{LOS-Sensitivität: Plasmaparameter}

\begin{itemize}
\item $ P_{ECRH} $ -Korrelation: höher besser
\item Optimale $ f_{rad} $ : 0,4-0,75
\item $ T_e \sim $ 1,3 keV ideal
\item Validiert hohe $ f_{rad} $ -Szenarien
\end{itemize}

\subsection{DAQ-Latenz}

\begin{itemize}
\item Laser-basierte Charakterisierung
\item Abtastzeiten: 0,8-3,2ms getestet
\item Minimum: 13,6ms erreicht
\item Totale Schleife: 150-400ms ausreichend
\end{itemize}

\end{document}
